\section{Troubleshooting and known limits}
\label{sec:troubleshooting}

Every refusal is a sentence naming the fault and its numbers; the session log
(\menu{File > Open log folder}) keeps it with everything else.

\begin{center}
\small
\begin{tabularx}{\textwidth}{@{}>{\raggedright\arraybackslash}p{0.34\textwidth}L@{}}
\toprule
Seen & What to do \\
\midrule
\app{Rack warning} on the status bar & Read its tooltip: a fault \oracle{} cannot mend (the
    old tagger wrapper) or hardware changed since the last session (measure the settle
    model and the bench table again) \\
\app{Run rejected} after Start scan & The gate's failing check is named
    (Section~\ref{sec:calibration}); nothing was written. A \emph{held} or \emph{settle}
    refusal points at the meter or the settling; repeat, and measure the settle time \\
A timing refused before starting & Shorten the gate, release or delay, or lengthen the
    dwell, as the sentence says \\
No triggers reach the time tagger & Check the PulseBlaster, its cable to the tagger and the
    trigger level; the scan holds until triggers return \\
A meter shows NaN & It stopped answering its poll: check its USB and that no other program
    holds it \\
The settings are refused at start-up & Fix the named line (on Windows, use single quotes
    around a folder, or forward slashes) \\
The acquisition process stopped & The run is closed on what it recorded; restart \oracle{} \\
A queue stopped & Its row says why; fix it and start the queue again, which takes up at the
    first scan not done \\
\bottomrule
\end{tabularx}
\end{center}

\paragraph{Known limits.}
\begin{itemize}
    \item Nothing has yet run on the beamline; treat each hardware path as unverified until
        bench-tested on the rack PC.
    \item SETTLE comes from bench steps up to 2\,kV; the beamline's settle model should be
        measured across the jumps scans make, including the largest range join.
    \item The settle gate (an AND unit on PulseBlaster channel 18) is not wired: settling
        bunches are dropped in software, to within the USB latency.
    \item The DAC's codes are 0.305\,V apart, so a 1\,V step carries up to $\pm$0.15\,V of
        rounding that no correction removes; a step of a whole number of codes has none.
    \item The DAC correction drifts (offsets by up to 0.1\,V a day): keep its validity short.
    \item The live view keeps every hit of the run in memory, 11 bytes a hit: about 40\,MB
        an hour at 1000 hits a second.
\end{itemize}
